\begin{lemma}[Homogeneous quadruples]
For every infinite $X\subseteq\mathbb N$ and every two-coloring of the
four-element subsets of $X$, there are an infinite $Y\subseteq X$ and a
color $b$ such that every four-element subset of $Y$ has color $b$.
\end{lemma}
\begin{proof}
Let $F=\mathsf{CardinalityFront}(4,X)$, as defined in
\noderef{CardinalityFront}.  By
\noderef{CardinalityFrontIsFront}, the hypothesis that $X$ is infinite
gives $\mathsf{Front}(F,X)$.  Define

\[
S=\{s\mid s\in F\text{ and }c(s)=1\}.
\]

The first conjunct in this set-builder definition shows explicitly that
$S\subseteq F$: if $s\in S$, then $s\in F$.  We may therefore apply
Nash--Williams \noderef{NashWilliams} to $F$, $X$, and $S$.  It gives a
family $F'$ and a set $Y$ such that

\[
\mathsf{Front}(F',Y),\qquad F'\subseteq F,\qquad
F'\subseteq S\ \text{or}\ F'\cap S=\emptyset.
\tag{1}
\]

We first identify the base of $F'$.  The infinite-base clause of the
explicit front definition \noderef{Front} says that $Y$ is infinite.
We claim that $\emptyset\notin F'$.  If $u\in F'$, then
$u\in F$ by $F'\subseteq F$, and the definition of
\noderef{CardinalityFront} gives $u.\mathop{\rm card}=4$.  In
particular, $u$ cannot be empty, whose cardinality is $0$.  Thus no
member of $F'$ is empty.  The base disjunction in the definition of
\noderef{Front} is consequently forced to take its nontrivial
alternative, and

\[
\mathsf{FrontBase}(F')=Y.
\tag{2}
\]

Here the base is the set of elements occurring in members of $F'$, by
the definition of \noderef{FrontBase}.  Equation (2) implies
$Y\subseteq X$: if $y\in Y$, then there is an $u\in F'$ with $y\in u$;
the inclusion $u\in F'\subseteq F$ and
\noderef{CardinalityFront} give $u\subseteq X$, so $y\in X$.

It remains to show that every four-element subset of $Y$ belongs to
$F'$.  Fix a finite set $s$ with $s.\mathop{\rm card}=4$ and
$s\subseteq Y$.  Since $s$ is nonempty, let $m=\max s$ and define

\[
T=\{y\in Y\mid m<y\},\qquad Z=s\cup T.
\]

Every element of $s$ is at most $m$, while every element of $T$ is
greater than $m$.  Thus the elements of $Z$ at most $m$ are exactly the
four elements of $s$; in increasing order, $s$ is the first four-element
prefix of $Z$.

We check carefully that $Z$ is an infinite subset of $Y$.  The inclusion
$s\subseteq Y$ is the hypothesis, and $T\subseteq Y$ by its definition,
so $Z\subseteq Y$.  For infinitude, apply
\noderef{InfiniteSetEnumeration} to the infinite set $Y$ and choose an
increasing enumeration $x\in\mathsf{IncSeq}$ with
$\operatorname{range}(x)=Y$.  By the order-embedding definition of
\noderef{IncSeq}, the values of $x$ are strictly increasing and
distinct.  Induction gives $x(q)\geq q$ for every $q$: this is clear at
$q=0$, and strict increase gives
$x(q+1)>x(q)\geq q$, hence $x(q+1)\geq q+1$.  Consequently every
$x(q)$ with $q\geq m+1$ belongs to $Y$ and is greater than $m$, so it
belongs to $T$.  These infinitely many distinct values show that $T$,
and therefore $Z$, is infinite.

Use the density clause of $\mathsf{Front}(F',Y)$ from
\noderef{Front} with the infinite set $Z\subseteq Y$.  We obtain
$t\in F'$ and $\mathsf{ProperPrefixSet}(t,Z)$.  Unpacking the definition
in \noderef{ProperPrefixSet}, choose $n\in Z$ such that

\[
\forall k,\qquad k\in t\quad\Longleftrightarrow\quad
k\in Z\text{ and }k<n.
\tag{3}
\]

Because $t\in F'\subseteq F$, \noderef{CardinalityFront} gives
$t.\mathop{\rm card}=4$.  We now prove $t=s$ using this fact and the
four-element hypothesis on $s$.  If $n\leq m$, then $n\in Z$ implies
$n\in s$, since every element of $T$ is greater than $m$.  Equation (3)
then shows that every element of $t$ is an element of $s$ strictly less
than $n$.  Thus $t\subseteq s\setminus\{n\}$, whose cardinality is
$3$, contradicting $t.\mathop{\rm card}=4$.  Hence $m<n$.

For every $k\in s$, the maximality of $m$ gives $k\leq m<n$, and
$k\in Z$; equation (3) therefore gives $k\in t$.  Thus $s\subseteq t$.
If some $k\in t$ were not in $s$, then (3) would put $k$ in $Z\setminus
s=T$ and give $m<k<n$.  The four distinct elements of $s$, together
with this $k$, would all belong to $t$, contradicting
$t.\mathop{\rm card}=4$.  Therefore $t\subseteq s$, and so $t=s$.
Since $t\in F'$, this proves

\[
s\in F'\qquad\text{for every }s\text{ with }s.\mathop{\rm card}=4
\text{ and }s\subseteq Y.
\tag{4}
\]

Finally consider the two alternatives in (1).  If $F'\subseteq S$, then
(4) gives $s\in S$ for every four-element $s\subseteq Y$.  The
definition of $S$ then gives $c(s)=1$, so take $b=1$.  If instead
$F'\cap S=\emptyset$, (4) gives $s\in F'$ and hence $s\notin S$ for
every such $s$.  Since $s\in F$ by
\noderef{CardinalityFront}, the definition of $S$ says that
$s\notin S$ means $c(s)\ne1$; as $c(s)$ is a Boolean, $c(s)=0$.  Take
$b=0$.  In both cases, (2) has already given $Y\subseteq X$, the front
infinite-base clause gave that $Y$ is infinite, and the required uniform
color conclusion follows.
\end{proof}
